\section{POWERS. COUNTABLE SETS}

\begin{enumerate}
\def\labelenumi{\arabic{enumi}.}
\tightlist
\item
  Two sets E, F are said to be equipotent if they can be put in one-to-one correspondence.
\end{enumerate}

Two sets, each equipotent to a third, are themselves equipotent.

If E and F are equipotent, then \(\mathfrak{P}(\mathbf{E})\) and \(\mathfrak{P}(\mathbf{F})\) are equipotent.

If E and F, \(E'\) and \(F'\) , \(E''\) and \(F''\) are respectively equipotent, then \(E \times E' \times E''\) and \(F \times F' \times F''\) are equipotent. This proposition extends to the product of any number of sets.

\begin{enumerate}
\def\labelenumi{\arabic{enumi}.}
\setcounter{enumi}{1}
\tightlist
\item
  Let X and Y be two generic subsets of a set E. The relation ``X and Y are equipotent'' is an equivalence relation on \(\mathfrak{P}(\mathbf{E})\) . The equivalence class (with respect to this relation) to which X belongs is called the power (*) of X, and the set of these classes (i.e., the quotient of \(\mathfrak{P}(\mathbf{E})\) by the above relation) is called the set of powers of subsets of E.
\end{enumerate}

If E and F are two distinct sets, the relation ``X and Y are equipotent'' between a subset X of E and a subset Y of F is expressed by saying that the power of X and the power of Y are equivalent. In this way we have a one-to-one correspondence between a subset of the set of powers of subsets of E and a subset of the set of powers of subsets of F.

\begin{enumerate}
\def\labelenumi{\arabic{enumi}.}
\setcounter{enumi}{2}
\tightlist
\item
  Let E, F be any two sets, which may or may not be distinct, and let a (resp. b) be an element of the set of powers of subsets of E (resp. F). Then a is said to be less than b, or b greater than a, if there exists a one-to-one mapping of a subset X ⊂ E with power a into a subset Y ⊂ F with power b. If also a and b are not equivalent powers, than a is said to be strictly less than b, or b strictly greater than a.
\end{enumerate}

If \(\mathfrak{a}\) and \(\mathfrak{b}\) are equivalent, then \(\mathfrak{a}\) is both greater and less than \(\mathfrak{b}\) . Conversely, it is a theorem that if \(\mathfrak{a}\) is both greater and less than \(\mathfrak{b}\) , then \(\mathfrak{a}\) and \(\mathfrak{b}\) are equivalent. It follows in particular that the set of powers of subsets of a set \(\mathbf{E}\) is ordered by the relation `` \(\mathfrak{a}\) is less than \(\mathfrak{b}\) ''; whenever we speak of this set as an ordered set, it is always the ordering defined by this relation that we mean.

Furthermore, using Zorn's lemma (and hence the axiom of choice), it is shown that the set of powers of subsets of a set E is well-ordered.

\begin{enumerate}
\def\labelenumi{\arabic{enumi}.}
\setcounter{enumi}{3}
\tightlist
\item
  The power of a set \(\mathbf{E}\) is strictly less than that of the set \(\mathfrak{P}(\mathbf{E})\) .
\end{enumerate}

If \(f\) is a mapping of a set \(\mathbf{E}\) into a set \(\mathbf{F}\) , then the power of the image \(f(\mathbf{X})\) of any subset \(\mathbf{X}\) of \(\mathbf{E}\) is less than the power of \(\mathbf{X}\) .

\begin{enumerate}
\def\labelenumi{\arabic{enumi}.}
\setcounter{enumi}{4}
\tightlist
\item
  Let \((\mathbf{X}_t)_{i \in I}\) be a family of subsets of a set \(E\) such that \(\mathbf{X}_t \cap \mathbf{X}_x = \emptyset\) whenever \(i \neq x\) . Let \((\mathbf{Y}_t)_{i \in I}\) be a family of subsets of a set \(F\) , indexed by the same set \(I\) and such that the power of \(\mathbf{Y}_t\) is less than that of \(\mathbf{X}_t\) for all \(i \in I\) . Then the power of the union \(\bigcup_{i \in J} \mathbf{Y}_t\) is less than that of \(\bigcup_{i \in J} \mathbf{X}_t\) for all subsets \(J\) of \(I\) .
\end{enumerate}

If also \(Y_{t} \cap Y_{x} = \emptyset\) whenever \(t \neq x\) , and if \(X_{t}\) and \(Y_{t}\) are equipotent for all \(t \in I\) , then \(\bigcup_{t \in J} X_{t}\) is equipotent to \(\bigcup_{t \in J} Y_{t}\) .

In particular, if F and E are the same set, we see that the power of the union of a set of pairwise disjoint subsets of E depends only on the powers of these subsets; it is called the sum of these powers (thus this function is defined for a family \((\mathfrak{a}_{\iota})\) of elements of the set of powers only if there exists a family \((\mathbf{X}_{\iota})\) of pairwise disjoint subsets of E such that \(X_{t}\) has power \(a_{t}\) for each index t).

If \((\mathbf{X}_i)_{i\in I}\) and \((\mathbf{Y}_i)_{i\in I}\) are families of subsets of E and F, respectively, indexed by the same set I such that \(\mathbf{X}_i\) and \(\mathbf{Y}_i\) are equipotent for all \(i\) , then the products \(\prod_{i}\mathbf{X}_{i}\) and \(\prod_{i}\mathbf{Y}_{i}\) are equipotent.

\begin{enumerate}
\def\labelenumi{\arabic{enumi}.}
\setcounter{enumi}{5}
\item
  The set N of natural integers may be considered as the set of powers of finite subsets of an infinite set. The order relation `` \(x \leqslant y\) '' on N is just the relation ordering this set of powers, and the sum of two natural integers is a function identical with the sum of two powers as defined above.
\item
  A set is said to be countable if it is equivalent to a subset of the set N of natural integers. Every finite set is therefore countable; if it has n elements, it is equipotent to the interval \([0, n - 1]\) of the set N. Every countable infinite set is equipotent to N; in particular, every infinite subset of N has the same power as N.
\end{enumerate}

If E is an infinite set, there exists a partition of E into countable infinite sets; in particular, every infinite set has a power greater than that of N.

If E is an infinite set, the sets \(E \times E\) and \(E \times N\) are both equipotent to E, and the set of finite subsets of E is equipotent to E. In particular, \(N \times N\) is a countable infinite set.

\begin{enumerate}
\def\labelenumi{\arabic{enumi}.}
\setcounter{enumi}{7}
\tightlist
\item
  A sequence of elements of a set E is by definition a family of elements of E, indexed by the set N or a subset of N. A sequence whose index set is N is therefore written \((x_{n})_{n\in\mathbf{N}}\) , or more simply \((x_{n})\) when there is no likelihood of confusion. If n denotes a generic integer, \(x_{n}\) is said to be the general term of the sequence, or the nth term. The latter terminology is also used when n is replaced by a particular integer. The set of elements of a sequence is countable.
\end{enumerate}

A sequence is said to be infinite or finite according as the index set is an infinite or finite subset of N. The set of elements of a finite sequence is finite.

Every subfamily of a sequence is again a sequence, called a subsequence of the given sequence. Every subsequence of a finite sequence is a finite sequence.

A family of elements whose index set is \(N \times N\) , or a subset of \(N \times N\) , is called a double sequence. A double sequence indexed by \(N \times N\) is written \((x_{m,n})\) , or more simply \((x_{mn})\) if there is no risk of confusion. Similarly for sequences with more than two indices.

Two sequences \((x_{n})\) , \((y_{n})\) are said to differ only in the order of their terms if there exists a permutation \(f\) of the index set such that \(y_{n} = x_{f(n)}\) for all \(n\) .

With any family of elements \((x_{i})_{i\in I}\) whose index set \(I\) is countably infinite we may associate an infinite sequence, as follows: there exists a bijection \(n\to f(n)\) of \(N\) onto \(I\) ; putting \(y_{n} = x_{f(n)}\) , the sequence \((y_{n})\) is said to be obtained by ranging the family \((x_{i})\) in the order defined by \(f\) . Thus the sequences corresponding to two distinct bijections of \(N\) onto \(I\) differ only in the order of their terms.

Operating in the same way when I is finite, we obtain a finite sequence associated with the family \((x_{i})\) .

\begin{enumerate}
\def\labelenumi{\arabic{enumi}.}
\setcounter{enumi}{8}
\tightlist
\item
  The union, intersection, and product of a family \((\mathbf{X}_t)_{t\in I}\) of subsets of a set \(E\) are said to be countable if \(I\) is a countable set, finite if \(I\) is finite.
\end{enumerate}

If I is countable, and if the power of \(X_{t}\) is less than a given infinite power a for all \(t \in I\) , then the power of the union \(\bigcup_{t} X_{t}\) is less than a. If also

at least one of the \(X_{i}\) has power a, then \(\bigcup_{i}X_{i}\) has power a. In particular, every countable union of sets of power a also has power a; every countable union of countable sets is a countable set.

